\section{Cosine \texorpdfstring{$\rone$}{R1}: Structural Failure}
\label{sec:cosine_failure}

During GRPO training, 256-token truncation reduces $\dsyc$ AUROC from 0.827 (offline CV) to approximately 0.60, so the optimizer works with a substantially weaker signal than the offline metrics suggest.
We test whether even this degraded discriminability translates into a viable reward through four configurations that form a diagnostic chain (Table~\ref{tab:cosine_chain}): each resolves the failure mode of its predecessor but exposes a deeper limitation.

\begin{table}[t]
\centering
\caption{Cosine $\rone$ failure chain on Qwen3-8B (LoRA rank$=$64, GRPO with $R = w_1 \cdot \rone + w_3 \cdot \rthree$).}
\label{tab:cosine_chain}
\resizebox{\columnwidth}{!}{%
\begin{tabular}{@{}lllll@{}}
\toprule
Config & $\rone$ Formula & $w_1$ & Outcome & Root Cause \\
\midrule
F1a & $\tanh(\mathbf{h}\cdot\dsyc/20)$ & 0.25 & Gate saturation & Norm mismatch \\
F1b & $\cos(\mathbf{h},\, \dsyc)$ & 0.25 & $\rone \approx 0$ & Gap too small \\
F1c & +recalibration & 0.25 & Mode collapse & Reward discontinuity \\
F1d & $5\cos(\mathbf{h},\, \dsyc)$ & 0.40 & Mode collapse & $\rthree$ saturation \\
\bottomrule
\end{tabular}}
\end{table}

The diagnostic chain begins with a norm mismatch.
F1a's dot-product formulation ($\rone = \tanh(\mathbf{h} \cdot \dsyc / 20)$) saturates because residual-stream norms exceed the direction norm by $100$--$1{,}200{\times}$, zeroing the gradient by step 955.
Cosine similarity (F1b) eliminates this norm dependence but reveals a second problem: the gap between valid and sycophantic cosines averages only 0.098, too narrow for meaningful reward contrast.
Recalibration (F1c, \S\ref{ssec:recalib}) widens the gap (offline AUROC rises to 0.910) but introduces a third: each recalibration event creates a reward discontinuity that destabilizes advantage estimates, collapsing entropy to 0.069 within 300 steps.
Amplification (F1d, $5{\times}$ scaling) lifts $\rone$ above the noise floor but surfaces the fundamental limitation: $\rthree$ saturates at 0.9998 by step 300, leaving the noisy $\rone$, which conflates warranted revision with unwarranted capitulation, as the sole gradient source until terminal collapse at step 2160.

Truncation further degrades the signal: under the 256-token maximum completion length used in GRPO, $\dsyc$ AUROC drops from 0.76 (full response) to 0.60, so the offline discriminability (0.827) substantially overestimates the signal available during training.
Checkpoint-level evaluation confirms immediate collapse: $\tof{=}0$ from step 500 onward with no stable phase.
The probe reward trap (C2, \S\ref{sec:probe_trap}) arises within GRPO using full-length responses and is unaffected by this truncation limitation; we verify that no evaluation response reaches the 256-token limit (Appendix~\ref{app:training}).
The four configurations converge on a single conclusion: $\dsyc$ encodes the \emph{direction} of belief change but not whether that change is \emph{warranted}.
Valid revision and sycophantic capitulation both shift activations toward the user's position along $\dsyc$, so the cosine projection assigns similar scores to both.
This is not a calibration artifact (recalibration improves discriminability yet rewards still fail) but a semantic limitation of any continuous $\dsyc$-based reward.
Section~\ref{sec:probe_trap} examines whether discretizing the signal via a binary probe can escape this limitation.
